
\documentclass[12pt]{article}
\usepackage[utf8]{inputenc}
\usepackage[T1]{fontenc}
\usepackage[portuguese]{babel}
\usepackage{amsmath}
\usepackage{siunitx}
\usepackage{graphicx}
\usepackage{geometry}
\usepackage{parskip}
\usepackage{enumitem}
\usepackage{tikz}
\usetikzlibrary{arrows.meta,calc,decorations.pathmorphing}

\geometry{a4paper, left=1cm, right=1.5cm, top=1cm, bottom=1.5cm}

\newcommand{\cabecalho}{
    \hfill
    \begin{minipage}[t]{0.7\textwidth}
        \raggedleft
        \textbf{Teste Final, Fundamentos de Física}
    \end{minipage}
    \vspace{0.35cm}
}

\newcommand{\pontos}[1]{\hfill\textbf{(#1 valores)}}

\setlength{\parindent}{0pt}
\setlist[enumerate]{itemsep=4pt, topsep=4pt}

\begin{document}

\cabecalho

\textbf{Nome:} \rule{10cm}{0.4pt}
\hfill
\textbf{Data:} \rule{2.5cm}{0.4pt}

\vspace{0.2cm}

\textbf{Indicação:} Sempre que necessário, considere $g = 10\,\text{m/s}^2$.
Antes de efetuar os cálculos, converta as grandezas para unidades do SI.

\vspace{0.25cm}

\textbf{Exercício 1}. Considere o vetor $\vec{v}=(v_x,v_y)$ de módulo $4$ e
orientação definida pelo ângulo $\theta=150^\circ$, medido a partir do semieixo positivo
de $x$.

\begin{enumerate}[label=\textbf{\alph*)}]
    \item Calcule as componentes de $\vec{v}$ e escreva o vetor usando os vetores
    unitários $\vec{i}$ e $\vec{j}$. \pontos{1,0}
    \item Calcule o vetor unitário com a mesma orientação do vetor $(v_x,-v_y)$. \pontos{1,0}
    \item Considere o vetor $\vec{w}=(3,\,3\sqrt{3})$. Determine $|\vec{w}|$ e o
    ângulo $\phi$, tais que $w_x=|\vec{w}|\cos\phi$ e
    $w_y=|\vec{w}|\sin\phi$. \pontos{1,0}
    \item Um vetor $\vec{u}$ tem módulo $6$ e forma um ângulo de $60^\circ$ com
    $\vec{v}$. Calcule o produto escalar $\vec{v}\cdot\vec{u}$ e explique o significado
    geométrico do produto escalar entre dois vetores. \pontos{1,0}
\end{enumerate}

\vspace{0.25cm}

\textbf{Exercício 2}. Uma moeda é lançada do solo na direção vertical, com
rapidez inicial de $10\,\text{m/s}$. Desprezando a resistência do ar, determine:

\begin{enumerate}[label=\textbf{\alph*)}]
    \item Qual é o instante em que a moeda atinge a altura máxima e qual é a sua velocidade
    nesse ponto? \pontos{1,5}
    \item Qual é a velocidade da moeda no instante em que atinge o solo? \pontos{1,0}
    \item A que distância horizontal do ponto de lançamento a moeda atinge o solo? \pontos{1,0}
\end{enumerate}

\vspace{0.25cm}

\textbf{Exercício 3}. Dois blocos estão ligados por uma corda ideal e sobem, sem atrito,
um plano inclinado que forma um ângulo de $30^\circ$ com a horizontal. Uma força de módulo
$140\,\text{N}$, paralela ao plano inclinado e dirigida para cima, puxa o bloco $m_1$.
As massas são:
\[
m_1=12\,\text{kg}, \qquad m_2=8000\,\text{g}.
\]

\begin{center}
	\begin{tikzpicture}[scale=0.95, >=Latex]
		\coordinate (A) at (0,0);
		\coordinate (B) at (8,0);
		\coordinate (C) at (8,4.62);
		\draw[thick] (A) -- (B) -- (C) -- cycle;
		\begin{scope}[shift={(2.8,1.62)}, rotate=30]
			\draw[fill=white, thick] (0,0) rectangle (1.6,0.95);
			\node at (0.8,0.48) {$m_2$};
			\draw[thick] (1.6,0.48) -- (3.2,0.48);
			\node[above] at (2.4,0.48) {$T$};
			\draw[fill=white, thick] (3.2,0) rectangle (4.8,0.95);
			\node at (4.0,0.48) {$m_1$};
			\draw[->, very thick] (4.8,0.48) -- (6.4,0.48) node[right] {$\vec{F}$};
		\end{scope}
		\draw (1.1,0) arc[start angle=0,end angle=30,radius=1.1];
		\node at (1.45,0.32) {$30^\circ$};
	\end{tikzpicture}
\end{center}

\begin{enumerate}[label=\textbf{\alph*)}]
	\item Desenhe o diagrama de corpo livre de cada bloco, indicando as forças paralelas
	e perpendiculares ao plano inclinado. \pontos{1,0}
	\item Calcule a aceleração do sistema. \pontos{1,5}
	\item Calcule a tensão $T$ na corda. \pontos{1,0}
	\item Considerando como sistema os dois blocos e a corda, indique as forças externas
	aplicadas ao sistema e as forças internas. Por qual princípio, na equação do sistema
	completo, as forças internas se cancelam? Explique. \pontos{1,0}
\end{enumerate}




\cabecalho

\textbf{Exercício 4}. Uma roldana gira inicialmente a $28\,\text{rev/s}$ e reduz
uniformemente a sua velocidade angular para $16\,\text{rev/s}$ durante um intervalo
de tempo total de $6,0\,\text{s}$.

\begin{enumerate}[label=\textbf{\alph*)}]
	\item Calcule a aceleração angular da roldana. \pontos{1,0}
	\item Calcule o número de revoluções completadas durante a primeira metade do
	intervalo de tempo total. \pontos{1,0}
	\item Calcule o ângulo total percorrido durante todo o intervalo de tempo. \pontos{1,0}
	\item Considere agora uma partícula que descreve uma circunferência com rapidez
	constante. Explique por que existe aceleração, apesar de a rapidez ser constante.
	Indique a direção e o sentido dessa aceleração. \pontos{1,0}
\end{enumerate}



\vspace{0.25cm}

\textbf{Exercício 5}. Uma esfera de massa $m=400\,\text{g}$ move-se horizontalmente
com rapidez inicial $v=15\,\text{m/s}$ e colide com um bloco de massa
$M=1600\,\text{g}$, inicialmente em repouso sobre uma superfície horizontal sem atrito.
O bloco está ligado a uma mola ideal de constante elástica $k=288\,\text{N/m}$ e
encontra-se inicialmente na posição de equilíbrio da mola, com $x=0$. A esfera fica
incrustada no bloco. Após o impacto, o conjunto bloco--esfera desloca-se para a direita,
alongando a mola.

\begin{center}
	\begin{tikzpicture}[scale=0.9, >=Latex]
		% Superfície e parede à esquerda
		\draw[thick] (-0.1,0) -- (11.0,0);
		\draw[thick] (0.5,0) -- (0.5,2.2);
		\foreach \y in {0.2,0.5,...,2.0}{
			\draw[thin] (0.5,\y) -- (0.2,\y+0.2);
		}
		% Mola abaixo da trajetória da esfera: ao mover-se para a direita, o bloco alonga a mola
		\draw[thick,decorate,decoration={zigzag,segment length=7pt,amplitude=4pt}]
		(0.5,0.30) -- (7.1,0.30);
		% Bloco
		\draw[fill=white, thick] (7.1,0) rectangle (8.6,1.1);
		\node at (7.85,0.62) {$M$};
		\node[above] at (3.8,0.52) {$k$};
		% Esfera incidente, acima da mola
		\draw[fill=white, thick] (4.9,0.82) circle (0.20);
		\node[below] at (4.9,1.57) {$m$};
		\draw[->, very thick] (5.15,0.82) -- (6.7,0.82) node[midway, above] {$\vec{v}$};
		\node at (5.8,-0.42) {};
	\end{tikzpicture}
\end{center}

\begin{enumerate}[label=\textbf{\alph*)}]
	\item Classifique a colisão. Indique as formas de energia envolvidas. Após o impacto,
	de que forma de energia para que forma de energia ocorre a transformação e porquê? \pontos{1,5}
	\item Calcule a rapidez do conjunto bloco--esfera imediatamente após o impacto. \pontos{1,5}
	\item Qual é a rapidez do conjunto no ponto de alongamento máximo da mola? \pontos{1,0}
	\item Calcule o alongamento máximo da mola e o trabalho realizado pela força elástica
	desde a posição de equilíbrio até ao ponto de alongamento máximo. \pontos{2,0}
\end{enumerate}


\end{document}
